\chapter{Making the generated projects build}
\label{ch:building}

Chapter~\ref{ch:whattobind} said what to bind. This chapter is about the other
half: getting the emitted projects to actually compile and link, which for a
real library means telling rosetta where the \emph{definitions} are, what
preprocessor state the headers expect, and which optimisation to build at.

\section{The two ways to supply definitions}

Your bound headers usually only \emph{declare} the API. The bodies live either
in a pre-built library or in source files. Rosetta has one field for each, and
they are independent --- use either, or both.

\begin{center}
\small
\begin{tabular}{@{}lL{10.2cm}@{}}
\toprule
\field{user\_lib} & Link the bindings against a pre-built \code{.so} /
  \code{.dylib} / \code{.a}. \\
\field{user\_sources} & Compile \code{.cpp} (or \code{.c}) files \emph{straight
  into} each binding target. \\
\bottomrule
\end{tabular}
\end{center}

\subsection{Linking a pre-built library}

\begin{lstlisting}[style=json]
"user_lib": {
  "name": "space",
  "dir":  "../space/bin",
  "link": "shared"
}
\end{lstlisting}

\begin{center}
\small
\begin{tabular}{@{}lcL{8.4cm}@{}}
\toprule
\textbf{Field} & \textbf{Req.} & \textbf{Meaning} \\
\midrule
\field{name} & \req & Library name --- the \code{space} in \code{libspace.dylib}. \\
\field{dir}  & \req & Directory holding it, relative to the manifest; used for
  \code{-L} and rpath. \\
\field{link} & & \code{"shared"} (default), \code{"static"}, or
  \code{"dynamic"} (an alias of shared). The \emph{preferred} form, with
  fallback to whichever is actually built. \\
\bottomrule
\end{tabular}
\end{center}

A bound library rarely stands alone. Give \field{user\_lib} an \textbf{array},
one object per library:

\begin{lstlisting}[style=json]
"user_include": ["../mylib/include", "/opt/foo/include"],
"user_lib": [
  { "name": "mylib", "dir": "../mylib/bin", "link": "shared" },
  { "name": "foo",   "dir": "/opt/foo/lib" },
  { "name": "bar",   "dir": "/opt/foo/lib", "link": "static" }
]
\end{lstlisting}

\begin{itemize}
  \item \textbf{Order is the link order} --- list a library before the ones it
    depends on, which is what static archives require.
  \item Each entry keeps its own \field{link}, so a shared library and a static
    dependency mix freely.
  \item Every distinct \field{dir} lands in the binding's \code{BUILD\_RPATH} and
    \code{INSTALL\_RPATH}, so the module loads without
    \code{LD\_LIBRARY\_PATH} / \code{DYLD\_LIBRARY\_PATH}.
  \item The \textbf{generator driver} links the same list --- it instantiates
    the bound types during the reflection walk, so it needs the dependencies'
    definitions too, not just yours.
  \item Headers of the dependencies go in \field{user\_include}; their
    configuration macros in \field{compile\_definitions}.
\end{itemize}

\begin{warn}
On \lang{wasm}, every entry is taken as the static archive \code{lib<name>.a}
built with the \textbf{same} emsdk. A native \code{.dylib} / \code{.so} cannot
enter a wasm module, so a dependency you cannot rebuild for emscripten rules
that target out.
\end{warn}

Flags that are not a \code{-L}/\code{-l} pair --- a framework, a
\code{--start-group}, a package's own link line --- go through a target's
\field{link\_options} (Chapter~\ref{ch:targets}); the two compose.

\subsection{Compiling sources in}

\begin{lstlisting}[style=json]
"user_sources": [
  "../src/widget.cpp",
  "../src/shape.cpp"
]
\end{lstlisting}

Always a list of paths, each relative to the manifest (or absolute); a single
string is accepted as a one-element list. Every compiled backend adds them via
\code{target\_sources(\ldots)}, so they build with the same include path and
flags as the generated binding.

Entries may be \textbf{shell globs}, expanded at generation time:

\begin{lstlisting}[style=json]
"user_sources": ["extern/arch/src/**/*.cxx"]
\end{lstlisting}

\code{*}, \code{?} and \code{[\ldots]} work within a path component, and a
component that is exactly \code{**} matches \textbf{zero or more directories}
--- so one pattern covers a whole source tree. Like the other wildcards,
\code{**} does not enter hidden directories. Matches are sorted for reproducible
output and the final list is de-duplicated, so mixing a literal path with a glob
that also covers it is safe. A pattern matching nothing warns and is skipped.

\begin{warn}
Globs expand when \code{rosetta\_gen} runs. \textbf{Re-run it after adding or
removing matching files} --- a new \code{.cpp} does not appear in the build by
itself.
\end{warn}

Entries may also be \textbf{C sources}. When any \code{.c} file is listed, the
generated \code{CMakeLists.txt} calls \code{enable\_language(C)} automatically,
which is what makes a library's vendored dependencies work:

\begin{lstlisting}[style=json]
"user_sources": [
  "./geogram/src/lib/geogram/mesh/*.cpp",
  "./geogram/src/lib/geogram/third_party/rply/*.c",
  "./geogram/src/lib/geogram/third_party/zlib/*.c"
]
\end{lstlisting}

The text-only backends compile nothing and ignore the field.

\section{Preprocessor definitions}

Most often a third-party library's configuration macros. Each entry is
\code{"NAME"} or \code{"NAME=VALUE"}:

\begin{lstlisting}[style=json]
"compile_definitions": [
  "GEOGRAM_USE_BUILTIN_DEPS",
  "GEOGRAM_WITH_HLBFGS"
]
\end{lstlisting}

They are emitted as \code{target\_compile\_definitions(\ldots\ PRIVATE \ldots)}
in \textbf{two} places, so the whole pipeline sees a consistent configuration:

\begin{itemize}
  \item the \textbf{generator driver} --- it includes the bound headers for the
    reflection walk, which must see the same preprocessor state the bindings
    will be built with;
  \item \textbf{every compiled binding target}, where they reach the bound
    headers and the \field{user\_sources} alike.
\end{itemize}

That first bullet is the one that matters. A macro that changes a class's member
set --- and plenty do --- would otherwise produce a binding for a different
class than the one you compile.

A single string is accepted as a one-element list.

\section{Build type and optimisation}
\label{sec:buildtype}

\begin{lstlisting}[style=json]
"build_type": "Release",
"optimization": "-O2"
\end{lstlisting}

Both optional and independent.

\field{build\_type} is one of \code{Debug}, \code{Release},
\code{RelWithDebInfo}, \code{MinSizeRel} (case-insensitive). It is emitted as a
\emph{default} inside
\code{if(NOT CMAKE\_BUILD\_TYPE AND NOT CMAKE\_CONFIGURATION\_TYPES)}, so
\code{-DCMAKE\_BUILD\_TYPE=\ldots} at configure time still wins and multi-config
generators (Visual Studio, Xcode) are left to pick a configuration at build
time.

\begin{note}
\textbf{Omitted means \code{Release}.} It used to mean ``emit nothing'', which
left CMake with no build type and hence no optimisation flags at all --- so the
documented two-line build produced an unoptimised module, around eight times
slower on a trivial pybind11 call. Pass \code{-DCMAKE\_BUILD\_TYPE=Debug}, or
set \code{"build\_type": "Debug"}, when you want an unoptimised build.
\end{note}

\field{optimization} is an explicit level: \code{-O0} through \code{-O3},
\code{-Os}, \code{-Oz}, \code{-Og} or \code{-Ofast} (the leading \code{-} may be
omitted). It is emitted as \code{add\_compile\_options(\ldots)} and
\code{add\_link\_options(\ldots)}, which land \emph{after} the build type's own
per-config flags --- so this \code{-O} overrides the level the build type
implies. \code{"build\_type": "Release", "optimization": "-O2"} builds
\code{-DNDEBUG} but at \code{-O2} instead of Release's \code{-O3}. The link
option matters for the wasm targets, where emscripten optimises at link time
too.

Both apply to the bindings and to any \field{user\_sources} compiled into them,
in every compiled backend.

\section{Headers the library's own build produces}
\label{sec:generatedheaders}

A header configured by the bound library's build system is simply absent when
rosetta compiles that library's sources without running its build. geogram's
\code{<geogram/version.h>} is the case in point --- configured by geogram's
CMake from \code{version.h.in}.

\begin{lstlisting}[style=json]
"generated_headers": [
  {
    "path": "geogram/version.h",
    "template": "./extern/geogram/src/lib/geogram/basic/version.h.in",
    "substitutions": {
      "VORPALINE_VERSION_MAJOR": "1",
      "VORPALINE_VERSION_MINOR": "10",
      "VORPALINE_VERSION_PATCH": "1",
      "VORPALINE_VERSION": "1.10.1-rosetta",
      "VORPALINE_BUILD_NUMBER": "",
      "VORPALINE_BUILD_DATE": "",
      "VORPALINE_SVN_REVISION": ""
    }
  }
]
\end{lstlisting}

\begin{center}
\small
\begin{tabular}{@{}lcL{8.4cm}@{}}
\toprule
\textbf{Field} & \textbf{Req.} & \textbf{Meaning} \\
\midrule
\field{path} & \req & The \emph{relative include path} --- exactly what the
  sources \code{\#include}. \\
\field{template} & one of & A file with \code{@KEY@} placeholders (CMake's
  \code{configure\_file} form), resolved against the manifest's directory. \\
\field{content} & one of & Literal lines, as an array of strings, for a header
  with no template. \\
\field{substitutions} & & \code{@KEY@} to value, applied to \field{template}. \\
\bottomrule
\end{tabular}
\end{center}

The file is written to \code{<bindings>/include/<path>} and that directory is
placed \textbf{first} on every generated target's include path --- ahead of the
library's own sources, so a stale checked-in copy of the same header cannot win
the lookup.

Substitution happens when the manifest is read, so a missing key is reported
against your manifest (\code{no substitution for @VORPALINE\_VERSION@}) rather
than compiling a literal \code{@KEY@} into the binding. Only the \code{@KEY@}
form is recognised, since \code{\$\{KEY\}} would collide with the shell and CMake
expansions such templates carry.

\begin{note}
A \code{-D} define is a different tool for a nearby job. geogram's
\code{GEOGRAM\_VERSION} is passed by its CMake on the command line, not written
into the header, so it belongs in \field{compile\_definitions}.
\end{note}

\section{Python wheels}
\label{sec:wheels}

The \lang{python} and \lang{nanobind} backends emit, alongside
\code{CMakeLists.txt}, everything needed to build a redistributable wheel:

\begin{center}
\small
\begin{tabular}{@{}lL{9.6cm}@{}}
\toprule
\code{pyproject.toml} & scikit-build-core build config. It drives the generated
  \code{CMakeLists.txt} itself, so the wheel and a plain \code{cmake --build}
  compile the same target the same way --- there is no second build description
  to keep in sync. \\
\addlinespace
\code{make\_wheel.py} & One-command builder: builds the wheel, then bundles
  external shared libraries and fixes the platform tag. Python rather than a
  shell script, so Linux, macOS and Windows are covered by one file. \\
\bottomrule
\end{tabular}
\end{center}

\begin{lstlisting}[style=sh]
cd bindings/nanobind
python make_wheel.py                    # -> dist/*.whl, dist/repaired/*.whl
python3.11 make_wheel.py                # a specific interpreter
python make_wheel.py --outdir /tmp/whl  # somewhere other than ./dist
python make_wheel.py --no-repair        # skip the bundling / retagging step
\end{lstlisting}

The wheel is built for whichever interpreter runs the script --- there is no
\code{--python} flag, just invoke the one you want (or pin it with the target's
\field{python} field, Section~\ref{sec:runtimepins}).

\field{version} is the distribution version stamped into the packaging
artifacts, PEP 440 style (\code{"1.2.0"}, \code{"0.3.0rc1"}). Default
\code{0.1.0}.

Two fields make packaging automatic for a project that always ships wheels.
They go \textbf{on the target}, beside the interpreter it is built for:

\begin{lstlisting}[style=json]
"targets": [
  { "lang": "python",   "name": "pygeom", "python": "3.11" },
  { "lang": "nanobind", "name": "nbgeom", "wheel": true, "wheel_dir": "./dist" }
]
\end{lstlisting}

\begin{center}
\small
\begin{tabular}{@{}lL{9.6cm}@{}}
\toprule
\field{wheel} & Build a wheel for \emph{this} target on every \code{--build},
  without passing \code{--wheel}. \\
\addlinespace
\field{wheel\_dir} & Where this target's wheels go instead of its own
  \code{dist/}. Implies \field{wheel}, and resolves against the manifest's
  directory. \\
\bottomrule
\end{tabular}
\end{center}

Only the two backends that emit a \code{make\_wheel.py} accept them:
\field{wheel} on a \lang{markdown} or \lang{node} target is an error, not a
no-op.

Per-target rather than top-level, and for the same reason \field{out\_dir} is
(Section~\ref{sec:outdir}): the two Python backends genuinely differ here.
\lang{nanobind} produces one \code{abi3} wheel covering every CPython 3.12+
where \lang{python} produces one per version, and \field{python}
(Section~\ref{sec:runtimepins}) already pins per target which interpreter each
is tagged for. Sharing one \field{wheel\_dir} between them is safe --- each
\code{make\_wheel.py} repairs only the wheels it just produced.

The \textbf{flags still win, in the direction of doing more}. \code{--wheel}
packages every \lang{python} / \lang{nanobind} target whatever the manifest
says, so a target spelling \code{"wheel": false} cannot turn off a run that
asked for it; \code{--wheel-dir DIR} overrides every target's own directory,
which is what makes ``collect them all in one place'' a single flag. With
neither set anywhere, \code{--build} only compiles the extension module.

\begin{warn}
\field{wheel} and \field{wheel\_dir} were \textbf{top-level} keys. They are
per-target now, and a top-level one is a load error naming the fix rather than
a key that silently does nothing --- a manifest that used to ship wheels must
not quietly stop.
\end{warn}

\begin{note}
On 3.12+, nanobind wheels are tagged \code{abi3}, covering every later CPython
with a single artifact.
\end{note}

\section{A realistic manifest}

Putting the last four chapters together --- this is roughly the shape a real
library's manifest takes:

\begin{lstlisting}[style=json]
{
  "//": "Bindings for the geom library.",

  "variables": [
    { "name": "SRC", "value": "./extern/geom" }
  ],

  "user_include": ["$SRC/include", "./extern/eigen-3.4.0"],
  "rosetta_include": "./extern/rosetta/include",
  "generator_name": "geom_gen",
  "module_name": "geom",

  "user_sources": ["$SRC/src/**/*.cpp"],
  "compile_definitions": ["GEOM_NO_ASSERT"],
  "build_type": "Release",
  "version": "1.2.0",

  "interop": ["eigen"],
  "sequences": [{ "type": "Eigen::VectorXd" }],

  "namespace": "geom",
  "header_dir": "geom",

  "targets": [
    { "lang": "python",   "name": "pygeom", "requires_python": ">=3.10" },
    { "lang": "nanobind", "name": "nbgeom", "out_dir": "./dist",
      "wheel": true, "wheel_dir": "./dist/wheels" },
    { "lang": "node",     "name": "jsgeom" },
    "typescript",
    "markdown"
  ],

  "classes": [
    { "header": "*.h", "exclude": ["fwd.h", "detail/*.h"] },
    { "header": "Mesh.h", "annotations": "Mesh.ann.json" }
  ],

  "functions": [
    { "name": "remesh", "header": "algorithms.h", "doc": "Isotropic remeshing." }
  ]
}
\end{lstlisting}
